\section{Exact-Quota Companion Processes}
\subsection{Exact CRT Quotas With Random Locations}
The random parent model fixes two harmful copy indices per parent. A separate statistical companion instead retains the exact number of accepted strikes supplied by a chosen CRT population and randomizes only their locations. This preserves the real count while removing the arithmetic targeting information, but it does not preserve two losses per 2-gap parent or the balanced global recurrence. Exact quotas create dependence within one layer, but they do not change the one-position survival scale when allocated uniformly on the conditioned eligible population. The head result also uses persistent availability and cross-layer mixing, while the square-window result uses blind placement and a quadratic eligible supply.
At filter $r$, let $U_r$ contain $N_r$ eligible values and let the CRT quota be $J_r$, with $0\le J_r\le N_r-2$. The exact-quota random parent model chooses one uniformly random size-$J_r$ subset of $U_r$ as its strike set. For a specified 2-gap whose two endpoints belong to $U_r$, both endpoints survive precisely when every strike is selected from the other $N_r-2$ values. Therefore
\begin{equation*}
\begin{aligned}
s_r
&=\frac{\binom{N_r-2}{J_r}}{\binom{N_r}{J_r}}
&&[\text{Uniform Exact-Quota Choice}]\\
&=\frac{(N_r-J_r)(N_r-J_r-1)}{N_r(N_r-1)}.
&&[\text{Factorial Simplification}]
\end{aligned}
\end{equation*}
Write the strike fraction as
\begin{equation*}
u_r:=\frac{J_r}{N_r}.
\end{equation*}
The exact formula gives
\begin{equation*}
\log s_r
=-2u_r+O\left(u_r^2+\frac{u_r}{N_r}\right).
\end{equation*}
This separates the exact finite quota from the cumulative condition needed for recurrence. Assume that along the conditioned chain to head $Q$,
\begin{equation*}
\begin{aligned}
\sum_{r < Q}u_r
&=\log\log Q+O(1),
&&[\text{CRT-Rate Cumulative Quota}]\\
\sum_{r < Q}
\left(u_r^2+\frac{u_r}{N_r}\right)
&=O(1).
&&[\text{Summable Finite-Population Error}]
\end{aligned}
\end{equation*}
The complete-period CRT benchmark $u_r=1/r$ satisfies these conditions. A different local CRT quota must be checked against them; preserving a numerical strike count alone does not make the conclusion automatic.
Multiplying the exact without-replacement factors gives
\begin{equation*}
\begin{aligned}
P_{\mathrm{quota}}(Q)
&=\prod_{r < Q}s_r
&&[\text{Survive Every Filter}]\\
&=\exp\left(\sum_{r < Q}\log s_r\right)
&&[\text{Product To Sum}]\\
&=\exp\left(-2\sum_{r < Q}u_r+O(1)\right)
&&[\text{Summable Error}]\\
&\asymp\frac{C}{(\log Q)^2}.
&&[\text{Cumulative Quota Condition}]
\end{aligned}
\end{equation*}
Thus the exact-quota companion has the same one-head survival order as the random parent model. It is not an independent Bernoulli filter inside one layer; it is a uniform shuffle conditioned on the exact CRT strike count.
Suppose the distinguished head candidate is eligible with conditional probability at least $b_0 > 0$, uniformly for all sufficiently large prime heads, and that this availability is compatible with the quota-survival experiment. Then
\begin{equation*}
\begin{aligned}
b_0P_{\mathrm{quota}}(Q)
&\le\Pr(H_Q)\le P_{\mathrm{quota}}(Q),
&&[\text{Uniform Availability Bounds}]\\
\Pr(H_Q)
&\asymp\frac{1}{(\log Q)^2}.
&&[\text{Quota Survival Asymptotic}]
\end{aligned}
\end{equation*}
Consequently,
\begin{equation*}
\begin{aligned}
\sum_{Q\text{ prime}}\Pr(H_Q)
&\asymp
\sum_{Q\text{ prime}}\frac{1}{(\log Q)^2}\\
&=\infty.
&&[\text{Prime Number Theorem}]
\end{aligned}
\end{equation*}
Under independence or an adequate cross-layer mixing condition, the second Borel-Cantelli lemma yields
\begin{equation*}
\Pr(H_Q\text{ occurs infinitely often})=1.
\qquad[\text{Q.E.D.}]
\end{equation*}
This is an almost-sure theorem, not a guarantee for every random realization. The set of realizations with only finitely many head hits has probability zero, but it is not logically empty.
For square-safe windows, assume $B(Q)\asymp C_0Q^2$ eligible starts and the same blind-placement empty-window premise used by the random parent model. Then
\begin{equation*}
\begin{aligned}
\lambda_{\mathrm{quota}}(Q)
&=B(Q)P_{\mathrm{quota}}(Q)
&&[\text{Expected Surviving Starts}]\\
&\asymp
C_0\frac{Q^2}{(\log Q)^2}
\longrightarrow\infty.
&&[\text{Quadratic Supply Dominates}]
\end{aligned}
\end{equation*}
If
\begin{equation*}
\Pr(W_Q\text{ is empty})
\le e^{-\lambda_{\mathrm{quota}}(Q)},
\end{equation*}
the empty probabilities are summable over prime $Q$. The first Borel-Cantelli lemma then gives only finitely many empty square windows almost surely. This eventual-window statement is stronger than the twin-prime-style target: an unbounded sequence of successful windows, or infinitely many head hits, is already sufficient for infinitely many distinct certificates.
For consecutive primes $p < q$, the real accepted-strike count in the next safe window is proved in \href{https://github.com/thiagomata/prime-numbers/blob/master/articles/chapter6/gap-dynamics.md#91-exact-accepted-strikes}{Gap Dynamics \S{}9.1} \cite{gapdynamics2026}:
\begin{equation*}
A(p,q)
=
\pi\left(\left\lfloor\frac{q^2-1}{p}\right\rfloor\right)
-\pi(p-1).
\end{equation*}
Using this local quota as $J_r=A(p,q)$ in the random-location companion is well defined, but its fractions $u_r=J_r/N_r$ must still satisfy the displayed cumulative conditions for the head proof above. The \href{https://github.com/thiagomata/prime-numbers/blob/master/articles/chapter6/gap-dynamics.md#52-exact-non-recursive-global-count}{complete-period count} \cite{gapdynamics2026} supplies the global density; neither exact count determines local placement in the real sieve.
\subsection{Biased Exact Quotas and the Logarithmic Skew Frontier}
The neutral exact-quota companion treats every eligible value symmetrically. We now keep the same quota $J_r$ while making 2-gap endpoints proportionally more likely to receive a harmful strike. This asks how much positional preference the quota can carry before the almost-sure head conclusion changes. We use a group-exchangeable allocation law, assume double strikes on one endpoint pair have quadratic order, and retain the availability, placement, and mixing premises of Section~3.5 and Section~7.1.
Let $E_r\subseteq U_r$ be the eligible values that are endpoints of locally relevant 2-gaps, and define
\begin{equation*}
x_r:=\frac{|E_r|}{N_r},
\qquad
u_r:=\frac{J_r}{N_r}.
\end{equation*}
Stipulate a group-exchangeable size-$J_r$ allocation law in which every endpoint has marginal inclusion probability $p_r^{E}$, every ordinary value has marginal inclusion probability $p_r^{O}$, and the endpoint preference ratio is
\begin{equation*}
\beta_r:=\frac{p_r^{E}}{p_r^{O}}\ge1.
\end{equation*}
We require that these marginals are feasible:
\begin{equation*}
0\le p_r^{O},p_r^{E}\le1.
\end{equation*}
This condition, together with the fixed-size law, is an assumption on the group-exchangeable quota construction rather than a consequence of the ratio $\beta_r$ alone.
The fixed quota forces the average marginal inclusion probability to equal $u_r$. Therefore
\begin{equation*}
\begin{aligned}
x_rp_r^{E}+(1-x_r)p_r^{O}
&=u_r
&&[\text{Exact Quota}]\\
p_r^{O}
&=\frac{u_r}{1+(\beta_r-1)x_r}
&&[\text{Substitution}]\\
p_r^{E}
&=\frac{\beta_ru_r}{1+(\beta_r-1)x_r}.
&&[\text{Simplification}]
\end{aligned}
\end{equation*}
Define the quota-normalized preference
\begin{equation*}
\kappa_r^{\mathrm{eff}}
:=
\frac{\beta_r}{1+(\beta_r-1)x_r}.
\end{equation*}
For one 2-gap, assume the probability that both endpoints are struck is $O((p_r^{E})^2)$. Its destruction fraction then satisfies
\begin{equation*}
\begin{aligned}
f_r
&=2p_r^{E}-\Pr(\text{both endpoints are struck})
&&[\text{Inclusion-Exclusion}]\\
&=2u_r\kappa_r^{\mathrm{eff}}
+O\left((u_r\kappa_r^{\mathrm{eff}})^2\right).
&&[\text{Normalized Preference}]
\end{aligned}
\end{equation*}
To use the displayed linear term cumulatively, we additionally require $p_r^E=u_r\kappa_r^{\mathrm{eff}}\to0$ and
\begin{equation*}
\sum_{r<\infty}\left(u_r\kappa_r^{\mathrm{eff}}\right)^2<\infty.
\end{equation*}
Without this error control, the per-filter expansion does not determine the asymptotic cumulative hazard.
The raw weight $\beta_r$ and the effective destruction skew are not identical: the exact quota must take probability away from ordinary values when it gives more probability to endpoints. For the complete-period eligible-sieve population, let $V_r$ be its accepted-value count and $T_r$ its 2-gap count. After filter $3$, 2-gaps have disjoint endpoints, so $|E_r|=2T_r$ and
\begin{equation*}
\begin{aligned}
x_r
&=\frac{2T_r}{V_r}\\
&=2\prod_{3\le p<r}\frac{p-2}{p-1}\\
&=\Theta\left(\frac1{\log r}\right).
&&[\text{Mertens-Type Product Estimate}]
\end{aligned}
\end{equation*}
The prior $\Theta((\log r)^{-2})$ density would be the density among raw integers, not among the eligible values $U_r$ used in this quota. For the illustrative sharper asymptotics
\begin{equation*}
x_r\sim\frac{C}{\log r},
\qquad
\beta_r\sim b\log r,
\end{equation*}
we instead obtain
\begin{equation*}
\kappa_r^{\mathrm{eff}}
=\frac{b}{1+bC}\log r\,(1+o(1)).
\end{equation*}
Thus a logarithmic raw preference changes its leading effective coefficient. A sparse relation $x_r=O((\log r)^{-2})$ may be imposed as an additional assumption for a different local target population, but it is not the complete-period eligible-sieve density.
For the phase theorem, measure skew by the realized effective factor
\begin{equation*}
\kappa_r
:=\frac{f_r}{2/r}.
\end{equation*}
This is the same quantity called $w_r$ in the general hazard analysis. Its cumulative survival is defined once $2\kappa_r < r$. Every regime below satisfies this inequality for all sufficiently large filters; the finite prefix is absorbed into a positive constant. Thus
\begin{equation*}
P_{\kappa}(Q)
=
\prod_{r < Q}\left(1-\frac{2\kappa_r}{r}\right)
=e^{-D_{\kappa}(Q)},
\end{equation*}
where
\begin{equation*}
D_{\kappa}(Q)
:=
\sum_{r < Q}-\log\left(1-\frac{2\kappa_r}{r}\right).
\end{equation*}
If $\kappa_r=\kappa < \infty$ is fixed, then
\begin{equation*}
\begin{aligned}
D_{\kappa}(Q)
&=2\kappa\log\log Q+O(1)
&&[\text{Prime Harmonic Sum}]\\
P_{\kappa}(Q)
&\asymp\frac{C_{\kappa}}{(\log Q)^{2\kappa}}.
&&[\text{Exponentiation}]
\end{aligned}
\end{equation*}
The sum of this probability over prime heads diverges for every finite $\kappa$. With persistent head availability and adequate cross-layer mixing,
\begin{equation*}
\text{every fixed finite proportional skew gives infinitely many head hits almost surely.}
\end{equation*}
Thus there is no finite constant-skew maximum.
The first transition appears when effective skew grows logarithmically. Set
\begin{equation*}
\kappa_r=1+c\log r,
\qquad c\ge0.
\end{equation*}
Then
\begin{equation*}
\begin{aligned}
D_c(Q)
&=2\log\log Q+2c\log Q+O(1)
&&[\text{Prime-Sum Asymptotics}]\\
P_c(Q)
&\asymp\frac{C_c}{Q^{2c}(\log Q)^2}.
&&[\text{Exponentiation}]
\end{aligned}
\end{equation*}
For prime heads, the occurrence series has the same convergence behavior as
\begin{equation*}
\int^\infty\frac{dx}{x^{2c}(\log x)^3}.
\end{equation*}
Therefore
\begin{equation*}
\begin{aligned}
c < \frac12
&\Longrightarrow
\text{infinitely many head hits almost surely, with mixing},\\
c\ge\frac12
&\Longrightarrow
\text{only finitely many head hits almost surely}.
\end{aligned}
\qquad[\text{Q.E.D.}]
\end{equation*}
In this explicitly normalized family, the equality $c=1/2$ is on the failure side because the remaining factor $(\log Q)^{-2}$ makes the boundary prime series converge. If $\beta_r$, rather than effective $\kappa_r$, is specified at its raw equality boundary, quota normalization changes lower-order terms and the cumulative series $D_{\kappa}(Q)$ must be evaluated directly.
The robust head-safe frontier is therefore
\begin{equation*}
\kappa_r
\le
1+\left(\frac12-\varepsilon\right)\log r
\quad\Longrightarrow\quad
\text{head recurrence almost surely, with mixing}.
\end{equation*}
For square windows, quadratic supply gives the larger robust frontier
\begin{equation*}
\kappa_r
\le
1+(1-\varepsilon)\log r
\quad\Longrightarrow\quad
\text{eventual square-window occupancy almost surely}.
\end{equation*}
Hence the intermediate range between approximately $(1/2)\log r$ and $\log r$ preserves square windows but not infinitely recurring head hits. The window statement is stronger than necessary: the twin-prime-style conclusion needs only infinitely many successful windows.
For an irregular skew schedule, the head criterion is
\begin{equation*}
\sum_{Q\text{ prime}}e^{-D_{\kappa}(Q)}=\infty,
\end{equation*}
together with persistent availability and adequate mixing. A pointwise skew percentage cannot replace this cumulative test.
